\chapter{Acute Bronchitis}

\section*{Definition}
Acute bronchitis is a self-limited inflammation of the tracheobronchial epithelium, characterized by acute cough (lasting less than 3 weeks) with or without sputum production. It is diagnosed clinically when you have an acute respiratory infection with cough as the predominant symptom, in the absence of pneumonia.

\section*{When to See a Doctor}

\begin{itemize}
 \item High or persistent fever, rapid breathing, confusion, or feeling markedly worse
 \item Cough producing blood-tinged sputum (coughing up blood)
 \item Shortness of breath at rest, blue lips, or low oxygen saturation
 \item Symptoms persisting beyond 3 weeks without improvement
 \item Cough that is worsening, lasts longer than expected, or keeps recurring, particularly with asthma, COPD, immune suppression, or other chronic disease
\end{itemize}

\section*{How It Develops}
Acute respiratory infection, usually viral, irritates the airways and can increase mucus and cough sensitivity. Temporary airway narrowing or wheezing may persist after the infection improves.

\section*{Causes}
\begin{itemize}
 \item \textbf{Viral (90\%):} Influenza A and B, rhinovirus, respiratory syncytial virus, adenovirus, coronavirus, parainfluenza
 \item \textbf{Bacterial (uncommon):} \textit{Bordetella pertussis} and other bacteria; pneumonia must be considered when systemic or examination findings are concerning
 \item \textbf{Environmental:} Tobacco smoke, air pollution, chemical irritants, dust
 \item \textbf{Other:} Pre-existing chronic bronchitis or asthma with superimposed infection
\end{itemize}

\section*{Related Symptoms}
\begin{itemize}
 \item Cough (dry or productive), sputum production, fever, fatigue, wheezing, chest tightness, sore throat
\end{itemize}
\textbf{Important:} Antibiotics are not routinely indicated for acute bronchitis in otherwise healthy adults, as the vast majority of cases are viral in origin.

\section*{Medical Evaluation}
\begin{itemize}
 \item Clinical diagnosis based on history and physical exam; chest imaging not routinely required
 \item Chest X-ray is considered when pneumonia is suspected, such as with abnormal vital signs, low oxygen, or focal lung findings; it is not routine for uncomplicated bronchitis
 \item Testing for influenza or SARS-CoV-2 may be useful when the result will change isolation, treatment, or contact advice
 \item Pertussis testing, usually PCR early in the illness, when the cough is paroxysmal or accompanied by whoop, vomiting, or relevant exposure
 \item Spirometry or bronchoprovocation testing if recurrent episodes suggest underlying asthma
\end{itemize}

\section*{Treatment Options}
\begin{itemize}
 \item Supportive care and, when appropriate, short-term symptom relief; evidence for many cough medicines is limited
 \item A bronchodilator may be considered when wheezing or underlying asthma or COPD is present, but it is not routine for uncomplicated bronchitis
 \item Antivirals may be prescribed for influenza or COVID-19 according to current guidance, risk factors, timing, and test or clinical assessment
 \item Antibiotics are not routinely indicated. They may be used for pertussis or another bacterial infection when confirmed or strongly suspected, or for pneumonia
 \item Cough suppressants (dextromethorphan, benzonatate) for sleep-disruptive cough; prescribe judiciously
\end{itemize}

\section*{Self-Care and Home Management}
\begin{itemize}
 \item Increased fluid intake (water, warm tea) to thin secretions and soothe irritated airways
 \item Honey (1--2 teaspoons) for cough suppression (avoid in children under 1 year)
 \item A clean humidifier may soothe irritation; avoid very hot steam because of burn risk
 \item NSAIDs or acetaminophen for fever and chest discomfort
 \item Avoid smoking and secondhand smoke exposure during recovery
\end{itemize}

\section*{References}
\begin{thebibliography}{99}
\bibitem{bronchitis_acute:ref1} Kinkade S, Long NA. ``Acute Bronchitis.'' Am Fam Physician. 2024;110(1):46--53.
\bibitem{bronchitis_acute:ref2} Smith MP, Lown M, Singh S, et al. ``Acute Cough Due to Acute Bronchitis in Adults.'' Chest. 2020;158(6):2299--2309.
\bibitem{bronchitis_acute:ref3} Harris AM, Hicks LA, Qaseem A. ``Appropriate Antibiotic Use for Acute Respiratory Tract Infection in Adults.'' Ann Intern Med. 2016;164(6):425--434.
\bibitem{bronchitis_acute:ref4} Irwin RS, French CL, Chang AB, et al. ``Classification of Cough as a Symptom in Adults and Children.'' Chest. 2018;153(1):196--209.
\bibitem{bronchitis_acute:ref5} Wenzel RP, Fowler AA. ``Clinical Practice: Acute Bronchitis.'' N Engl J Med. 2006;355(20):2125--2130.
\end{thebibliography}
